%% Write your paper content here.
%% Each \section, figure, table, etc. goes in this file.
%% Keep all figures inside the figures/ subfolder next to this file.
\section*{MLIR dialects and the 'linalg' Dialect}
Multilayered Intermediate Representation (MLIR) gets its name by being a framework that contains many dialects at different levels of abstraction. Thess Dialects are the mechanism by which to engage with and extend the MLIR ecosystem. They allow for defining new attributes, operations, and types. 

In other words, MLIR achieves a transformation that takes IR as input and produces modified IR as output transformation, these transformations are called “passes”.

The following is the exploration of MLIR dialect “linalg”, a dialect designed for linear algebra computations. We will explore linalg by adding two vectors in parallel. Starting with a high level IR, performing multiple passes, till we lower to LLVM IR and then to an executable.
\\
\\
\begin{figure}[H]
\centering
\begin{tikzpicture}[
    node distance=1.5cm,
    box/.style={rectangle, draw, 
                minimum width=2.5cm, minimum height=0.8cm}
]

\node[box] (node1) {High Level IR};
\node[box] (node2) [below=of node1] {Dialect lowering Passes};
\node[box] (node4) [below=of node2] {LLVM IR};
\node[box] (node5) [below=of node4] {Executable};



\draw[->] (node1.south) -- ++(0,-0.75) -| (node2.north);
\draw[->] (node2.south) -- ++(0,-0.75) -| (node4.north);
\draw[->] (node4.south) -- ++(0,-0.75) -| (node5.north);

\end{tikzpicture}
\caption{Overview of the compilation flow}
\end{figure}

\section*{High-level IR}
We start by defining a parallel vector addition in linalg.generic, which is an operation from the ‘linalg’ dialect that allows us to describe almost any type of tensor operations without having to worry about things like loop structure; this is the highest level of abstraction, you can find this in "vec\_add.mlir".

\begin{minted}{llvm}
-linalg-to-parallel-loops
module {
  func.func @parallel_vector_add(
    %arg0: memref<1024xf32>, 
    %arg1: memref<1024xf32>, 
    %arg2: memref<1024xf32>
    ) {
    %c0 = arith.constant 0 : index
    %c1024 = arith.constant 1024 : index
    %c1 = arith.constant 1 : index
    scf.parallel (%arg3) = (%c0) to (%c1024) step (%c1) {
      %0 = memref.load %arg0[%arg3] : memref<1024xf32>
      %1 = memref.load %arg1[%arg3] : memref<1024xf32>
      %2 = arith.addf %0, %1 : f32
      memref.store %2, %arg2[%arg3] : memref<1024xf32>
      scf.reduce 
    }
    return
  }
}
\end{minted}


\section*{Dialect Lowering Passes}
Now that we have defined the parallel vector addition using linalg.generic, we need to lower this IR to LLVM IR and then get our executable, to do so we will multiple MLIR passes from different dialect including one from linalg, this is where MLIR shines and allows you to reuse components of different passes to optimize for your specific targets, if we were working with a single pass it would be really hard to maintain and wouldn't be reusable. The passes are as follow:

\begin{center}
\begin{tabular}{||c c p{7cm}||} 
 \hline
 Pass & Dialect in -> Dialect out & Description \\ [0.5ex] 
 \hline\hline
 --convert-linalg-to-parallel-loops & linalg->scf & Converts abstract operations (linalg.generic) into explicit parallel loops (scf.parallel). \\ 
 \hline
 --convert-scf-to-openmp & scf->omp & Converts abstract operations (linalg.generic) into explicit parallel loops (scf.parallel). \\
 \hline
 --convert-openmp-to-llvm & omp->llvm & Converts parallel loops (scf.parallel) into OpenMP constructs (omp.parallel, omp.wsloop). \\
 \hline
 --convert-scf-to-cf & scf->cf & Converts OpenMP ops into LLVM dialect ops and runtime library calls. \\
 \hline
 --convert-arith-to-llvm & arith->llvm & Converts structured control flow (for/while loops) into basic blocks and branches. \\ 
 \hline
 --convert-func-to-llvm & func->llvm & Converts arithmetic operations (arith.addf) into LLVM arithmetic (llvm.fadd). \\ 
 \hline
 --finalize-memref-to-llvm & memref->llvm & Converts function signatures and calls to LLVM function format. \\ 
 \hline
 --reconcile-unrealized-casts & (cleanup) & Removes temporary type conversion artifacts from previous passes. \\ 
 \hline
\end{tabular}
\end{center}

\section*{LLVM IR and executable}

After running these MLIR passes we can now translate the output IR into LLVM-IR, this is possible since MLIR has an LLVM dialect, 

From here, we have an LLVM file that can be compiled by the clang compiler, we just add a C file to serve as our entry point and populate the arrays, we can now compile using the following command:

\begin{minted}{bash}
clang -fopenmp main.c vec_add.ll -o vec_add
\end{minted}

\section*{Moving Forward}

We have now explored how we can leverage an existing MLIR dialect to make a multi-stage compilation for parallel computing. From here I see two different ways in which we can proceed with a research project for MLIR.

\subsection*{Building an MLIR Dialect for an HPC framework}
There are lots of built-in MLIR dialects that are great for many different applications, the idea of this project would be to find a specific HPC pattern that could benefit from an MLIR transformation and building an MLIR dialect for it.

\subsection*{Identify optimization opportunity for MLIR Dialect}
MLIR is great at abstraction, but mapping those abstractions to specialized HPC hardware has lots of room for improvement. The idea of this project is to identify a specific optimization that a existing MLIR dialect handles poorly and provide an improvement.


